\documentclass[a4paper]{esannV2}
\usepackage[T1]{fontenc}
\usepackage{lmodern}
\usepackage[utf8]{inputenc}
\usepackage{graphicx,amsmath,amssymb,array,booktabs,url}
% Official class retained unchanged. Correct its legacy offsets to current rules.
\voffset 0cm \hoffset 0cm
\setlength{\oddsidemargin}{18.6mm}\setlength{\evensidemargin}{18.6mm}
\setlength{\topmargin}{27.6mm}\setlength{\headheight}{0mm}\setlength{\headsep}{0mm}
\setlength{\textwidth}{122mm}\setlength{\textheight}{191mm}
\pdfpagewidth=210mm \pdfpageheight=297mm
\newcommand{\F}{F_1}
\newcommand{\BA}{\mathrm{BA}}
\clubpenalty=10000
\begin{document}
\title{When Disagreement Labels Do Not Identify\protect\\F1 Rankings}
\author{Roberto Fern\'andez-Barrios, Iker Pastor-L\'opez,\\
Asier Gonz\'alez-Santocildes and Pablo Garc\'ia Bringas
\vspace{.3cm}\\
Faculty of Engineering, University of Deusto\\
Avda. de las Universidades, 24, 48007 Bilbao, Spain
}
\maketitle
\begin{abstract}
When two binary classifiers agree, their labels cancel from an accuracy
comparison but can still change an F1 ranking. We specialize partial
identification to an explicit audit: after labeling all disagreements,
when can one choose a model, and when must agreements be queried?
The criterion certifies either an exact ranking or a choice within an F1
tolerance of the better model for every compatible completion. Across 3,200
frozen pairs, retaining the shared completion tightens bounds and a 0.01 tolerance
reduces query costs substantially. Validation-selected comparators retain
material ambiguity in 6 of 36 cells. These are finite-batch guarantees,
not population error estimates.
\end{abstract}

\section{What can disagreement labels decide?}
Suppose two fixed classifiers disagree on three of ten items. One predicts
positive on all three, the other negative; only one of these labels is positive.
They agree positive on three other items and negative on four. Even knowing
that exactly four of the ten labels are positive does not settle their F1
ranking. Allocating the three remaining positives to negative agreements gives
F1 scores $(1/5,0)$; allocating them to positive agreements gives $(4/5,6/7)$.
The same disagreement labels and prevalence support opposite winners.

We ask when disagreement labels suffice to choose between fixed models.
Our specialization yields a sharp sign criterion, an operational regret
certificate, and retrospective evidence on information and labeling costs.
The key is to compare models under the \emph{same} completion of unknown labels.

Active F-measure estimation and expected-loss comparison optimize acquisition
\cite{sawade2010,sawade2012}. Active Testing estimates risk without bias
\cite{kossen2021}; Bayesian Assessment uses practical-equivalence regions
\cite{ji2021}; CODA selects models using consensus priors and accuracy regret
\cite{kay2025}. Our guarantee instead covers every compatible completion.
F1 batch dependence \cite{lipton2014} and missing-label metric bounds
\cite{dervovic2025} are established; our specialization builds on TDC and Polo.

\section{An explicit audit criterion}
Fix binary predictions $q,c\in\{0,1\}^n$ and correct unknown labels.
Let $a,b$ count disagreements $(q,c)=(1,0),(0,1)$;
$n_{00},n_{11}$ count agreements. After all $D=a+b$ disagreement labels are
revealed, let $u,v$ be their respective positive counts. Unknown positive
agreement counts are $s,t$. Write
\[
 P=u+v+s+t,\quad Q=a+n_{11},\quad C=b+n_{11},\quad d=u-v.
\]
$P$ is the true positive count, while $Q,C$ are observable predicted counts.
Positive-class F1 is $2\mathrm{TP}/(P+\text{predicted positives})$, with
$0/0=0$. The paired difference is
\begin{equation}\label{eq:gap}
 \Delta(s,t)=\F(q)-\F(c)=\frac{2(t+u)}{P+Q}-\frac{2(t+v)}{P+C}.
\end{equation}
By contrast, accuracy difference $(2d+b-a)/n$ is already determined;
balanced accuracy difference $[nd+(b-a)P]/[2P(n-P)]$ is determined if
$0<P<n$ is known, and undefined at $P=0,n$.

\noindent\textbf{Sharp sign criterion (specialization).}
Set $e=d+b-a$, $K=d(u+v+n_{11})+ub-va$. Then $\Delta$ has the sign of
\begin{equation}\label{eq:affine}
 S(s,t)=K+ds+et.
\end{equation}
Without prevalence information, agreement queries leave an integer rectangle
$s\in[s_-,s_+]$, $t\in[t_-,t_+]$, initially $[0,n_{00}]\times[0,n_{11}]$.
Its exact numerator bounds are
\begin{align}\label{eq:bounds}
 S_-&=K+\min(ds_-,ds_+)+\min(et_-,et_+),\nonumber\\
 S_+&=K+\max(ds_-,ds_+)+\max(et_-,et_+).
\end{align}
Both strict rankings are possible iff $S_-<0<S_+$. A strict winner is
identified when $S_->0$ or $S_+<0$; an exact tie when $S_-=S_+=0$.
Remaining cases permit a tie and one winner.

\noindent\textit{Proof.} Cross multiplication gives
$\Delta(P+Q)(P+C)/2=dP+(C-Q)t+uC-vQ=K+ds+et$.
The denominator is positive unless $P=0$ and a predictor has no positives;
then both F1 scores and $S$ are zero. Every integer count in the rectangle
is attainable. An affine function reaches its extrema at its corners,
proving the sign claims.\hfill$\square$

\noindent\textbf{Equal-count corollary.}
If $Q=C=R>0$, then $\Delta=2d/(P+R)$: disagreement labels identify the sign
without knowing $P$. If $R=0$, both scores vanish. Magnitude can remain unknown.

\noindent\textbf{Exact magnitude bounds.}
For fixed $P>0$, the difference is affine in $t$:
\begin{equation}\label{eq:fixed}
 f_P(t)=\frac{2u}{P+Q}-\frac{2v}{P+C}
          +\frac{2(C-Q)t}{(P+Q)(P+C)}.
\end{equation}
The feasible endpoints are
$\ell_P=\max(t_-,P-u-v-s_+)$ and
$h_P=\min(t_+,P-u-v-s_-)$.
If $\ell_P\le h_P$, evaluating $f_P$ at both endpoints gives its exact hull;
otherwise that $P$ is infeasible. At $P=0$ both scores are zero. With unknown
$P$, enumerate feasible integer counts $P$ and take the smallest and largest
of these endpoints, yielding the exact interval hull $[L,U]$ in $O(n)$ time.
This includes $P=n$. The sign proof does not place unknown-$P$
magnitude extrema at rectangle corners.

\noindent\textbf{A choice with bounded regret.}
For absolute F1 tolerance $\varepsilon\ge0$, the worst losses from selecting
$q$ and $c$ are $R_q=\max(0,-L)$ and $R_c=\max(0,U)$.
Choose the smaller, breaking ties in favor of $q$, and stop iff
\begin{equation}\label{eq:stop}
 \min(R_q,R_c)\le\varepsilon
 \quad\Longleftrightarrow\quad L\ge-\varepsilon\ \text{or}\ U\le\varepsilon.
\end{equation}
For every compatible completion, choosing $q$ loses at most $\max(0,-L)$
relative to the better model; the analogous bound for $c$ is $\max(0,U)$.
Thus (\ref{eq:stop}) guarantees an $\varepsilon$-best choice on this batch.
It does not require $[L,U]\subseteq[-\varepsilon,\varepsilon]$.
At $\varepsilon=0$ it allows a weak winner before distinguishing a strict win
from an exact tie. Failure means $L<-\varepsilon$ and $U>\varepsilon$: opposite compatible
advantages exceeding the tolerance. The original material-reversal definition
$L\le-\varepsilon$, $U\ge\varepsilon$ differs only at the boundaries; the
regret criterion is the decision interpretation of this condition.

\medskip\noindent\textbf{Relation to general identification frameworks.}
TDC \cite{tdc}, Theorem~1, gives sharp certificates for fixed additive-loss
contrasts over unrestricted completions. With $P>0$ fixed,
\begin{equation}\label{eq:additive}
 \Delta=\sum_i y_i\left[\frac{2q_i}{P+Q}-\frac{2c_i}{P+C}\right].
\end{equation}
The cardinality constraint $\sum_i y_i=P$ must also be imposed; TDC's
balanced-accuracy construction (\S4.5) already does so and enumerates
prevalence counts. Its unrestricted itemwise lower bound alone need not be
sharp under this constraint. Equations (\ref{eq:affine})--(\ref{eq:fixed})
are explicit specializations, not a new general certificate theorem.

Polo et al.\ \cite{polo2024}, Eq.~(1.1) and \S3.2, also cover fixed-$P$
F1. Set $X$ uniform on batch indices; let $Z$ distinguish audited items and
group unaudited items. Observations and $P$ fix the $(Y,Z)$ marginal. Use $n$
times the summand in (\ref{eq:additive}) as the paired objective. Remaining
fractional labels lie in $[0,1]$ with integer sum. Every vertex is integral:
two fractional entries admit a fixed-sum perturbation, while one cannot sum
to an integer. Thus linear extrema equal deterministic-completion extrema.
We supply the explicit audit solution and its empirical decision consequences.

The separate interval is
$[\min\F(q)-\max\F(c),\allowbreak\max\F(q)-\min\F(c)]$. Every extremum uses the
\emph{same} feasible completions and prevalence constraints; their marginal
optimizers may nevertheless be incompatible. Fixed-$P$ marginals are affine
in $t$, so endpoint enumeration is exact. This isolates discarded dependence,
not superiority to joint TDC or Polo.

\section{Retrospective evaluation}
\textbf{Frozen cases and units.} We reuse TDC's prediction artifact, release
1.1.6 (distinct from preprint v1), without refitting models. Each of 36 cells
contains 500 targets, a validation-selected quantum-kernel predictor, and
115 classical configurations. Twenty cells cover BRFSS diabetes and ACS
FoodStamps (five subsampling seeds, SVC and Gaussian-process classification);
16 cover eight EMBER, ToN-IoT and UNSW-NB15 scenarios with both classifiers.
Deduplicating identical classical predictions within cells yields 3,200 pairs.
The criterion applies to arbitrary fixed binary predictions.

The eight-cell pilot (801 pairs: seed 42, BRFSS, ACS and both EMBER scenarios,
both classifiers) was fixed before reading labels in this study and met its
continuation criteria. The 28 reserved cells (2,399 pairs) are an analytical
holdout of existing evidence. Tolerance, separate-bound and validation-comparator analyses
are retrospective extensions after review, not prospective confirmation.

Fractions average equally across cells; costs average orders, pairs within
cells, then eligible cells. Shared targets and overlapping subsamples preclude
independent-observation inference; summaries are descriptive. Known $P$ is
an exact information oracle. Initial intervals retain $0<P<n$, the original
balanced-accuracy domain; query cohorts and stops allow $P=0,n$.

\textbf{Acquisition and comparators.} We retain three agreement orders:
uniform without replacement; alternating shuffled agreement strata; and
decreasing $|d|$ or $|e|$ with shuffled ties. Seeds 0--31 share base permutations
across policies and both stopping objectives. Orders use predictions and
revealed disagreement labels only. Queries shrink the count intervals. Magnitude ordering minimizes numerator
width at each budget, not necessarily stopping cost.
All 307,200 trajectories are replayed for the original strict-sign/exact-tie
objective and for $\varepsilon=0.01$. Full labels verify realized regret;
they do not select pairs, orders or stopping decisions.

We also compare one classical predictor per cell: the closest in Hamming
distance (original-order tie break), and the maximum validation
balanced-accuracy candidate (lexicographic configuration tie break).
Recovered validation scores cover all 115 candidates per cell; choices were
fixed before the new target analysis. This retrospective rule follows TDC's
objective, not validation F1. Ten cells have validation ties: alternative
tie choices preserve 6/36 and all extra costs, shifting mean totals by at most
one label. Hamming proximity is not a deployment rule.

\textbf{Mechanism.} At known $P>0$, interval width is exactly
$2|C-Q|(h_P-\ell_P)/[(P+Q)(P+C)]$: predicted-count imbalance and freedom to
allocate positives between agreement groups jointly drive ambiguity.
Without $P$, numerator width is $|d|n_{00}+|e|n_{11}$ initially.
All 58 equal-count pairs in 25 cells identify the sign, as predicted,
although their cell-mean gap width remains 0.094. Disagreement count alone
therefore does not describe decision-relevant uncertainty.

\textbf{Information after disagreements.} The original mean material-reversal
fractions are 21.93\% in the pilot and 20.54\% in the reserved cells; with
exact $P$ they are 15.22\% and 10.52\%. These count pairs admitting opposite
compatible winners with margins at least 0.01, not observed mistakes or a
universal frequency. Table~\ref{tab:intervals} translates ambiguity into decisions. Tolerance
percentages complement material reversals here: no pair lies on the relevant
0.01 boundaries. Separate bounds certify far fewer choices.

\begin{table}[!ht]
\centering
\begin{tabular}{llrrr}
\toprule
$P$ & Interval & Width & Exact (\%) & $\varepsilon$ (\%)\\
\midrule
Unknown & Paired & 0.173 & 68.44 & 79.46 \\
Unknown & Separate & 1.585 & 6.53 & 6.57 \\
Exact & Paired & 0.034 & 76.18 & 89.48 \\
Exact & Separate & 1.005 & 17.86 & 18.80 
\\
\bottomrule
\end{tabular}
\caption{Initial bounds in the 28 reserved cells. Equal-cell means of interval
width and certified decisions. Exact means strict winner or exact tie;
$\varepsilon=0.01$ means bounded regret. ``Exact'' $P$ is oracle information.}
\label{tab:intervals}
\end{table}

Allowing $P=0,n$ changes twelve paired gap intervals but no material-reversal
classification. Excluding the 14 constant classical predictors changes the
all-cell material mean only from 20.85\% to 20.87\%. Selecting the nearest
comparator leaves material ambiguity in \textbf{3/36} cells; validation selection leaves
\textbf{6/36} (both have 11/36 strict reversals). These one-pair-per-cell
designs answer a different question from the all-configuration pool.

\textbf{Label costs.} Table~\ref{tab:costs} uses common cohorts for both goals,
including pairs already decided. Under magnitude ordering, replacing the exact goal by tolerance reduces total
cost from 121.16 to 74.59 labels across all pairs and from 92.57 to 40.29
under validation selection, versus 500 for full labeling. Nearest-comparator
costs do not describe the entire pool.

\begin{table}[!ht]
\centering
\setlength{\tabcolsep}{3.5pt}
\begin{tabular}{llrrrrr}
\toprule
 & & Initial & \multicolumn{2}{c}{Additional} & \multicolumn{2}{c}{Total}\\
Cohort & Policy & $D$ & Exact & $\varepsilon$ & Exact & $\varepsilon$\\
\midrule
All pairs & Uniform & 54.5 & 88.21 & 35.58 & 142.70 & 90.07 \\
All pairs & Alternating & 54.5 & 80.64 & 31.29 & 135.13 & 85.78 \\
All pairs & Magnitude & 54.5 & 66.67 & 20.10 & 121.16 & 74.59 \\
Validation & Uniform & 30.7 & 89.07 & 25.56 & 119.79 & 56.28 \\
Validation & Alternating & 30.7 & 77.58 & 16.34 & 108.30 & 47.07 \\
Validation & Magnitude & 30.7 & 61.84 & 9.57 & 92.57 & 40.29 \\
Nearest & Uniform & 14.1 & 114.04 & 1.84 & 128.18 & 15.98 \\
Nearest & Alternating & 14.1 & 96.27 & 1.72 & 110.41 & 15.86 \\
Nearest & Magnitude & 14.1 & 80.31 & 1.04 & 94.45 & 15.18 
\\
\bottomrule
\end{tabular}
\caption{Mean labels on common cohorts, 36 cells, 32 orders per policy.
All pairs: 3,200 comparisons; validation and nearest: one per cell.
Additional labels query agreements. Every batch contains 500 items.}
\label{tab:costs}
\end{table}

Conditioning instead on initial \emph{exact} indecision gives 1,132 pairs in
34 cells, with mean initial cost 48.61. Extra exact/tolerance costs are
249.71/92.22 (uniform), 221.26/81.21 (alternating), and 185.27/53.50 (magnitude).
The 736 initially \emph{tolerance}-undecided pairs occupy 32 cells; their
initial cost is 57.39 and tolerance extras are 152.00, 134.30 and 87.37.
These conditional denominators must not be interchanged. For example, the
all-pair magnitude tolerance total is 82.05 under equal \emph{pair} weights,
versus the reported 74.59 under equal \emph{cell} weights.

Checks cover 7,506 rational hulls and 2,100 sequential/future-label cases;
original exact stops are unchanged. In 15,466 of 307,200 replays, the certified choice is the realized-worse model;
maximum regret is 0.00927, within tolerance. These are dependent replay counts.
The all-completion guarantee follows from (\ref{eq:stop}).

\section{Use and limitations}
The criterion tells an auditor whether to choose now or query agreements.
Retaining shared completions and matching the stopping goal to the decision
tolerance substantially reduce labeling in these cases.

This incremental specialization assumes fixed predictions and correct labels;
it gives no population, noisy-label, macro-F1, multiclass or family-wise guarantee.
Compatible extremes need not be probable; benchmark reuse limits breadth.
Observed costs are not optimal-acquisition lower bounds. Disagreement-first
labeling is imposed; separate pairwise costs cannot be summed across the family.
Statistical methods can trade assumptions or error probability for fewer labels.
Code, frozen derived inputs and results are archived at
\url{https://doi.org/10.5281/zenodo.23642981}.

\noindent\textbf{Declarations.}
R.F.-B. performed the research and writing; I.P.-L., A.G.-S. and P.G.B.
provided scientific guidance. No specific funding was received. OpenAI Codex
assisted with literature, design, mathematical checks, code, analysis and
writing; Claude assisted with independent review. All authors approved the
content and accept responsibility.
\begin{footnotesize}
\begin{thebibliography}{9}
\bibitem{sawade2010} C. Sawade, N. Landwehr, and T. Scheffer.
Active Estimation of F-Measures. \emph{NeurIPS}, 23, 2010.
\bibitem{sawade2012} C. Sawade, N. Landwehr, and T. Scheffer.
Active Comparison of Prediction Models. \emph{NeurIPS}, 25, 2012.
\bibitem{kossen2021} J. Kossen, S. Farquhar, Y. Gal, and T. Rainforth.
Active Testing: Sample-Efficient Model Evaluation. \emph{ICML}, PMLR 139:5753--5763, 2021.
\bibitem{ji2021} D. Ji, R. L. Logan IV, P. Smyth, and M. Steyvers.
Active Bayesian Assessment of Black-Box Classifiers. \emph{AAAI}, 35:7935--7944, 2021.
\bibitem{kay2025} J. Kay, G. Van Horn, S. Maji, D. Sheldon, and S. Beery.
Consensus-Driven Active Model Selection. \emph{ICCV}, 4594--4604, 2025.
\bibitem{lipton2014} Z. C. Lipton, C. Elkan, and B. Narayanaswamy.
Thresholding Classifiers to Maximize F1 Score. arXiv:1402.1892, 2014.
\bibitem{dervovic2025} D. Dervovic and M. Cashmore.
Model Evaluation in the Dark: Robust Classifier Metrics with Missing Labels.
\emph{AISTATS}, PMLR 258:1909--1917, 2025.
\bibitem{tdc} R. Fern\'andez-Barrios, I. Pastor-L\'opez, A. Gonz\'alez-Santocildes,
and P. Garc\'ia Bringas. Sharp Target-Domain Certificates for Quantum-Kernel
Advantage under Distribution Shift. \emph{arXiv:2610.00080v1}, 2026.
\url{https://arxiv.org/abs/2610.00080v1}.\bibitem{polo2024} F. M. Polo, S. Maity, M. Yurochkin, M. Banerjee, and Y. Sun.
Weak Supervision Performance Evaluation via Partial Identification.
\emph{NeurIPS}, 37, 2024. arXiv:2312.04601.
\end{thebibliography}
\end{footnotesize}
\end{document}
